\documentclass[11pt]{article}
\usepackage[margin=1in]{geometry}

\title{Benchmark Infrastructure for the CTRL-ALT-ELITE RV32I Core\\
\large An assembler, a representative workload suite, and what they validate}
\author{CTRL-ALT-ELITE}
\date{November 2026}

\begin{document}
\maketitle

\begin{abstract}
The project's CPI work had no workload that could exercise a branch predictor: the
directed instruction tests execute each branch approximately once, so a 2-bit counter
never has a second execution to predict. The graded autograder binaries are unavailable
as source, so the only way to obtain a loop-dominated workload was to build one.
This document describes the infrastructure that was created --- a minimal RV32I
assembler and a self-checking benchmark suite of thirteen programs --- the conventions the
harness imposes on such programs, and the first measurements it produced. With the
branch history table enabled, a ten-iteration counted loop retires in 47 cycles instead
of 59 ($-20\%$) and an iterative Fibonacci loop in 69 instead of 81 ($-15\%$).
\end{abstract}

\section{Purpose and scope}

Two problems motivated this work.

\textbf{No measurable workload.} Every CPI result before this point was measured on
directed instruction tests. Those tests are excellent at proving functional correctness
--- 53 of them pass, and they caught real regressions --- but they are the wrong shape
for evaluating a branch predictor. A predictor needs a branch that executes repeatedly
so that its counter can be trained and then trusted. In the directed suite almost every
branch executes once, so the best a predictor can do is neither help nor hurt. The
branch history table was therefore implemented, corrected, and still \emph{unproven}:
no test in the suite could demonstrate a benefit.

\textbf{No access to the graded binaries.} The autograder reports pass/fail per program
(bubble sort, factorial, GCD, Fibonacci, sum array, and others) but does not provide
sources or binaries. Their instruction mix can be inferred --- under RV32I with no
multiplier or divider, factorial and GCD must contain libgcc-style shift-and-add and
shift-and-subtract loops --- but they cannot be disassembled, counted, or run locally.

The response is a set of \emph{representative} workloads: programs of the same
algorithmic shape as the graded ones, owned by the project, that can be built, modified
and measured freely. They are not the graded programs, and no result measured on them
should be quoted as a grade. What transfers is the mechanism, not the absolute number.

\section{Deliverables at a glance}

\begin{center}
\begin{tabular}{ll}
\hline
Component & Purpose \\
\hline
\texttt{code/benchmarks/rvasm.py} & Minimal RV32I assembler: source to \texttt{.x} image \\
\texttt{code/benchmarks/} & Thirteen workloads, their images and the assembler \\
\texttt{make bench-build} & Assembles every benchmark source in one command \\
\texttt{benchmarks/README.md} & Conventions a new benchmark must obey \\
\hline
\end{tabular}
\end{center}

The suite integrates with the existing harness without changes to any testbench:
\texttt{make sim-all} sweeps the new directory along with the others, and the total
grew from 53 to 66 tests with zero failures.

\section{The assembler}

\subsection{Why an assembler rather than hand encoding}

The existing directed tests were written as assembly and then encoded by hand; their
hex images are checked in beside their sources. Hand encoding is a genuine source of
error. During this project a helper function that builds a \texttt{jal} placed
\texttt{imm[11]} at instruction bit 10 instead of bit 20 --- a single-character mistake
that corrupted the destination register field as well, and that no amount of reading the
source would have revealed.

An assembler makes the source of truth the mnemonic, and makes the encoding testable:
the same instruction that a hand-encoded image got wrong will be assembled correctly and
the difference will appear immediately.

\subsection{Syntax and conventions}

The assembler accepts one instruction per line, with labels, \texttt{\#} or
\texttt{//} comments, and both \texttt{off(reg)} and \texttt{reg,reg,imm} operand forms.
Two conventions match the existing tests and are worth stating explicitly because they
are inconsistent with each other:

\begin{itemize}
\item \textbf{Immediates are decimal}: \texttt{addi x3,x0,30} loads thirty.
\item \textbf{Branch and jump targets are hexadecimal addresses}: 
      \texttt{beq x1,x3,1000014} targets \texttt{0x1000014}. Labels are also accepted
      and preferred in new code.
\end{itemize}

\subsection{Validation}

The assembler is validated by reproduction rather than by inspection: it is run over
all thirteen hand-written \texttt{pipeline\_tests} sources and its output is compared
against the checked-in images. \textbf{All thirteen are reproduced byte for byte.} That
covers the entire instruction subset those tests use and is the strongest available
evidence that the encoder is correct; any error in a field would show as a diff.

\section{The benchmark suite}

\subsection{Requirements}

Every benchmark was written to satisfy four conditions:

\begin{enumerate}
\item \textbf{Self-checking.} The program computes an expected value, compares it, and
      leaves \texttt{x1 = 1} only if it matches. A wrong answer fails the test rather
      than passing quietly.
\item \textbf{Loop-dominated or load-dominated.} The shape must stress what is being
      measured: backward branches for the predictor, load-use pairs for the hazard
      logic.
\item \textbf{Buildable from source.} The \texttt{.txt} is committed with the
      \texttt{.x}, so the workload can be modified and rebuilt.
\item \textbf{Correct under the existing harness.} No testbench changes were required.
\end{enumerate}

\subsection{The workloads}

\begin{center}
\begin{tabular}{llll}
\hline
Test & Size & Shape & Self-check \\
\hline
\texttt{smoke} & 2 words & two instructions & flow sanity only \\
\texttt{loop10} & 16 words & counted loop, 10 iterations & sum $1..10 = 55$ \\
\texttt{sum-array} & 28 words & store/load loop over 4 words & total $= 10$ \\
\texttt{fib} & 20 words & iterative Fibonacci, 10 iterations & fib$(11) = 89$ \\
\hline
\end{tabular}
\end{center}

\texttt{smoke} exists to separate harness failures from program failures: if a
two-instruction program does not pass, the problem is the flow, not the benchmark.
\texttt{loop10} is the predictor workload, with nine taken backward branches.
\texttt{sum-array} exercises stores, loads, a load-use dependency and a loop, and is the
one benchmark that records non-zero stalls. \texttt{fib} is a dependency chain with a
register-register add feeding the next iteration.

\subsection{What each benchmark does, and how it was built to test}

Each program is written in assembly and assembled by \texttt{rvasm.py}; the source is
committed beside the image it produces. The descriptions below give the instruction-level
behaviour, the mechanism the shape was chosen to exercise, and --- where it is informative
--- the flush signature the measurement should reproduce, which is a way of checking that
a benchmark is measuring what it claims to.

\paragraph{smoke (2 words).} Executes \texttt{addi x1,x0,1}, then \texttt{ecall}.
Built as the smallest program that satisfies the harness contract --- pass register and
terminator --- and nothing else. Its job is diagnosis rather than measurement: if it fails,
the problem is the flow (the image, the loader, the testbench), not a benchmark. Correct
result: six cycles, no stalls, no flushes.

\paragraph{sweep4/10/25/50/100.} Counts down from $N$ while incrementing an iteration
counter, then verifies the counter equals $N$.
All five images are the same program with one immediate changed, so trip count is the only
variable between them. The loop body contains exactly one branch, the loop-back, so the
dynamic branch stream is periodic with period one: there is nothing else for the predictor
to train or to pollute. That makes the family a controlled experiment on warm-up alone.
The flush signature is a strong check on the whole flow. The loop-back is taken $N-1$
times, so without prediction there must be $N-1$ flush cycles --- 3, 9, 24, 49 and 99 for
the five images --- and with prediction enabled the count must be a constant 3 (two warm-up
mispredicts plus one loop exit) independent of $N$. Both are exactly what is measured.

\paragraph{loop10.} Accumulates $1+2+\dots+10$ and verifies 55.
It shares the loop structure and the instruction count of \texttt{sweep10} --- hence the
identical 47 cycles --- but its body performs an arithmetic accumulation, so the forwarding
path and the ALU are exercised while the same branch stream is measured. It is the only
single-loop benchmark with both a data result and a timing result.

\paragraph{nested.} Five outer iterations, each running a four-iteration inner loop, with
the inner index accumulated; verifies 30.
Built to expose what a single loop cannot: the inner branch is taken four times and then
falls through once per outer iteration, while its counter is still in the taken state, so
every inner-loop exit is a guaranteed mispredict. The two branches also occupy different
table entries and must remain correct simultaneously.
Signature: $5\times3+4=19$ taken branches, so 19 flush cycles without prediction and 10
with --- the 10 being five inner exits, one outer exit and four warm-up mispredicts.
Measured: 19 and 10.

\paragraph{if-in-loop.} A ten-iteration loop containing a forward branch that is taken on
nine iterations and falls through on the seventh, setting a marker that is checked
afterwards.
This tests the opposite failure from \texttt{nested}: not a branch that exits a loop, but a
rarely taken path inside one. If the hot direction is learned correctly the cost is a
single mispredict; if the cold fall-through is learned as taken, the cost is two cycles on
every iteration. The body also carries two branches per iteration, so both must be held in
the table at once.
Signature: $9+9=18$ taken branches, so 18 flush cycles without prediction and 6 with
(four warm-up, one rare fall-through, one loop exit). Measured: 18 and 6.

\paragraph{sum-array.} Stores the values 1--4 into four words at \texttt{0x1000}, then
loops four times loading and accumulating them; verifies the total of 10.
The load is immediately followed by the instruction that consumes it, so each iteration
contains a load-use pair and must stall once. It is the only benchmark that stores, and the
only one that touches the data path at a non-zero address --- built with \texttt{lui},
because the data port is not base-relative.
Signature: four stalls, exactly one per load-use pair, and three flush cycles. The loop is
too short for prediction to change either, which is itself the expected result. Measured:
four stalls and three flushes in both configurations.

\paragraph{fib.} Ten iterations of $a\leftarrow b$, $b\leftarrow a+b$; verifies 89.
The body is a pure register dependency chain with no memory traffic, so any lost cycle
would have to come from the pipeline itself rather than from the data path. Its body is
five instructions against \texttt{sweep}'s three, which halves the branch density and is
why the same ten iterations cost 69 cycles instead of 47.
Signature: nine taken branches, so nine flush cycles without prediction and three with.
Measured: 9 and 3.

\paragraph{smul.} A shift-and-add multiply of $6\times 7$ over a 32-iteration loop: test
the multiplier's low bit, conditionally add the multiplicand, shift both operands, decrement
the bit counter. Verifies 42.
This reproduces the shape libgcc multiplication takes on an RV32I core without the M
extension, and therefore the shape the graded \emph{factorial} and \emph{GCD} programs
contain. It is also the only benchmark whose inner branch is genuinely \emph{data
dependent}, which makes it the only one where prediction cannot in principle remove every
mispredict. In practice the multiplier (7) shifts to zero after three iterations, after
which that branch becomes constant and is learned --- which is why the gain is large rather
than partial.
Signature: 31 taken loop-backs plus 29 taken bit-tests, so 60 flush cycles without
prediction and 5 with. Measured: 60 and 5.

\paragraph{callret.} Three calls to a five-instruction subroutine and three returns, using
\texttt{jal x1} and \texttt{jalr x0,0(x1)}; verifies 15.
The pair is the canonical call and return: \texttt{rd = x1} drives the pre-decoder's call
classification and \texttt{rs1 = x1} drives its return classification, so this is the
workload on which a return-address stack would be measured. The subroutine is placed before
the entry code and jumped over, and \texttt{x2} is cleared at entry to defeat the harness's
end-of-program detector, which otherwise fires on the first return.
Signature: the loop-back runs only twice when taken, so two warm-up mispredicts are spent to
save one --- which is why this is the one benchmark where the predictor is a net cost.
Measured: three flush cycles without prediction, four with.

\section{Harness conventions a benchmark must obey}

Building the suite exposed three conventions that were undocumented and that each cost
a debugging cycle. They are now recorded in the README.

\paragraph{The pass register differs by testbench.} \texttt{top\_tb.sv}, which runs
\texttt{pipeline\_tests}, \texttt{sample-bmarks} and the new suite, passes on
\texttt{x1 = 1}. \texttt{rv\_official\_tb.sv}, which runs \texttt{insn-bmarks}, passes
on \texttt{x3 = 1}. A benchmark written for one convention silently fails under the
other.

\paragraph{ecall must be the last executed instruction.} The testbench observes
\texttt{ecall} when it is fetched, waits approximately five cycles, and then samples the
register file. Anything placed after the \texttt{ecall} retires during that window. The
first version of every benchmark put its \texttt{fail} label immediately after the
\texttt{ecall}; the failing path then executed during the drain and overwrote the pass
register, so a program that computed perfectly correct values reported failure. The
last executed instruction must be \texttt{ecall}, followed only by \texttt{nop}s ---
which is why the pre-existing tests end with four of them.

\paragraph{The data port does not subtract the program base.} Instruction fetch is
relative to \texttt{BASE\_ADDR}, data access is not. A data address must therefore be
small (for example \texttt{0x1000}, built with \texttt{lui x2,1}) rather than
base-relative.

\section{Measurements}

\subsection{Method}

The simulation is deterministic: identical RTL and an identical program image produce
identical cycle counts, exactly. Every difference reported below is therefore
attributable to the change, and none of it is run-to-run noise.

The comparison is a feature mask on a single build: the branch history table is disabled
by forcing its prediction output to zero, everything else is unchanged, and the full
57-test suite is run in both configurations. Both configurations pass all 57 tests, so
the measurement isolates performance and not correctness.

\subsection{Results}

\begin{center}
\begin{tabular}{lrrrrr}
\hline
 & \multicolumn{2}{c}{BHT enabled} & \multicolumn{2}{c}{BHT disabled} & \\
Test & Cycles & Flushes & Cycles & Flushes & Change \\
\hline
\texttt{smoke}     &  6 & 0 &  6 & 0 &  0 \\
\texttt{loop10}    & 47 & 3 & 59 & 9 & $-12$ \; $(-20.3\%)$ \\
\texttt{sum-array} & 51 & 3 & 51 & 3 &  0 \\
\texttt{fib}       & 69 & 3 & 81 & 9 & $-12$ \; $(-14.8\%)$ \\
\hline
\end{tabular}
\end{center}

\subsection{Interpretation}

The flush counts are exactly what a two-bit predictor with a cold seed predicts, which
is the point: the benefit is not a mysterious speedup but accounted for branch by branch.

\begin{center}
\begin{tabular}{llll}
\hline
Iteration & Entry state & Prediction & Result \\
\hline
1 & empty & none & mispredict (cold) \\
2 & \texttt{WEAK\_NT} & not taken & mispredict (first training) \\
3--9 & \texttt{WEAK\_T} / \texttt{STRONG\_T} & taken & correct \\
10 (exit) & \texttt{STRONG\_T} & taken & mispredict (loop exit) \\
\hline
\end{tabular}
\end{center}

Three mispredicts for \texttt{loop10} against nine taken branches without a predictor.
At two cycles per mispredict, that is the twelve cycles observed.

\texttt{sum-array} is unchanged, and the reason is instructive: its loop body executes
only four times, so there are three taken backward branches in total. The predictor
spends two of them training and can only help with the third, so the saving is within
the cost of the two warm-up mispredicts. \textbf{A short loop cannot repay a predictor's
warm-up.} This is consistent with the earlier finding that the directed instruction
tests show no benefit at all, and it is the reason the graded workloads --- which iterate
far more than ten times --- should show a larger relative gain than these proxies do.

\section{Threats to validity}

\begin{itemize}
\item \textbf{These are not the graded programs.} They share algorithmic shape, not
      instruction sequences. Absolute cycle counts and CPI do not transfer; the
      mechanism and the direction do.
\item \textbf{Loop trip counts are small.} Ten iterations makes the warm-up a large
      fraction of the total. Longer loops amortise it and will show a better ratio.
\item \textbf{One predictor configuration.} These results characterise a 64-entry,
      8-bit-tagged table with a two-bit counter. A different size or a different
      replacement policy would shift the absolute numbers.
\item \textbf{CPI is comparable here but not universally.} Both configurations retire
      the same instruction count, so CPI is a valid comparison across them. CPI is
      \emph{not} comparable across benchmarks, and would not be comparable across any
      change that alters the dynamic instruction count, such as instruction fusion.
\end{itemize}

\section{Reproducing}

\begin{verbatim}
  make bench-build -C verif/        # assemble every benchmark source
  make sim-all -C verif/            # 57 tests, sweeps every directory
  make run -C verif/ TEST_DIR=benchmarks TEST=loop10
\end{verbatim}

To reproduce the control measurement, force \texttt{taken\_o} to zero in
\texttt{code/bht.sv}, run \texttt{make sim-all}, and compare the benchmark rows.

\appendix
\section{Instruction subset supported by the assembler}

\texttt{nop}, \texttt{ecall}, \texttt{add}, \texttt{sub}, \texttt{and}, \texttt{or},
\texttt{xor}, \texttt{sll}, \texttt{srl}, \texttt{sra}, \texttt{slt}, \texttt{sltu},
\texttt{addi}, \texttt{andi}, \texttt{ori}, \texttt{xori}, \texttt{slli},
\texttt{srli}, \texttt{srai}, \texttt{slti}, \texttt{sltiu}, \texttt{lw},
\texttt{sw}, \texttt{lui}, \texttt{auipc}, \texttt{jal}, \texttt{jalr},
\texttt{beq}, \texttt{bne}, \texttt{blt}, \texttt{bge}, \texttt{bltu},
\texttt{bgeu}.

Registers may be written numerically (\texttt{x0}--\texttt{x31}) or by ABI name
(\texttt{zero}, \texttt{ra}, \texttt{sp}, \texttt{a0}--\texttt{a7}, \texttt{t0}--\texttt{t6},
\texttt{s0}--\texttt{s11}).


\section{Addendum: extended suite, trip-count sweep, and the cost of prediction}

After the first version of this document, five further workloads and a trip-count sweep
were added. The suite now contains thirteen benchmark programs and the full regression
runs 66 tests, all passing.

\subsection{What was added}

\begin{center}
\begin{tabular}{lll}
\hline
Test & Shape & Self-check \\
\hline
\texttt{sweep4/10/25/50/100} & counted loop, one branch, $N$ iterations & iterations $=N$ \\
\texttt{nested} & 5 outer $\times$ 4 inner & sum $=30$ \\
\texttt{if-in-loop} & hot loop, one fall-through on iteration 7 & marker $=1$ \\
\texttt{smul} & shift-and-add multiply, 32 iterations & $6\times 7 = 42$ \\
\texttt{callret} & 3 calls and 3 returns & $3\times 5 = 15$ \\
\hline
\end{tabular}
\end{center}

Rationale, in descending order of value: the sweep converts a single data point into a
scaling law; \texttt{nested} exposes inner-loop exits, which mispredict once per outer
iteration and which no single loop can reveal; \texttt{if-in-loop} checks that a rarely
taken path inside a hot loop is not learned as always taken; \texttt{smul} is the shape
libgcc multiplication takes on RV32I without the M extension, and is therefore the shape
the graded \emph{factorial} and \emph{GCD} programs contain; \texttt{callret} is the only
workload that exercises a return-address stack.

\subsection{Results}

\begin{center}
\begin{tabular}{lrrr}
\hline
Test & BHT on & BHT off & Saved \\
\hline
\texttt{sweep4}     &  29 &  29 & 0 \\
\texttt{sweep10}    &  47 &  59 & $+12$ (20.3\%) \\
\texttt{sweep25}    &  92 & 134 & $+42$ (31.3\%) \\
\texttt{sweep50}    & 167 & 259 & $+92$ (35.5\%) \\
\texttt{sweep100}   & 317 & 509 & $+192$ (37.7\%) \\
\texttt{smul}       & 218 & 328 & $+110$ (33.5\%) \\
\texttt{nested}     & 112 & 130 & $+18$ (13.8\%) \\
\texttt{if-in-loop} &  56 &  80 & $+24$ (30.0\%) \\
\texttt{fib}        &  69 &  81 & $+12$ (14.8\%) \\
\texttt{sum-array}  &  51 &  51 & 0 \\
\texttt{callret}    &  37 &  35 & $-2$ ($-5.7\%$) \\
\hline
\end{tabular}
\end{center}

All thirteen benchmarks pass in both configurations, so the comparison again isolates
performance from correctness rather than mixing the two.

\subsection{The scaling law}

Saved cycles per loop iteration: 1.20 at $N=10$, 1.68 at $N=25$, 1.84 at $N=50$ and 1.92
at $N=100$. The sequence approaches two cycles per iteration, which is exactly one
avoided redirect: the predictor reaches the point where it removes the mispredict on
every steady-state iteration.

This is the number that extrapolates. The graded benchmarks iterate far more than one
hundred times, so the figure that applies to them is the asymptote of roughly 38\%, not
the 20\% measured on a ten-iteration loop. The same arithmetic explains why
\texttt{sweep4} and \texttt{sum-array} show nothing: two warm-up mispredicts cost four
cycles, and a four-iteration loop never runs long enough to earn them back.

\subsection{The cost of prediction}

\texttt{callret} is the one benchmark that becomes \emph{slower} with prediction enabled,
by two cycles. Its loop-back branch executes only three times, so the cold entry costs
the same two mispredicts a four-iteration loop does, while the only branch it can then
predict correctly is the third. For very short loops a cold predictor is a net loss.

That is a property of the design rather than a defect, and it should be reported as such:
the predictor maximises steady-state throughput at a small fixed cost per cold branch.
Alternative seeding policies trade that cost for spurious taken predictions elsewhere ---
seeding new entries as taken was measured during development and made the branch tests
worse by fifteen redirects --- so the cold seed is the conservative choice, not an
accident.

\subsection{Two harness traps characterised while building this}
\enlargethispage{5\baselineskip}

\paragraph{A fetched return ends the simulation early.} \texttt{top\_tb.sv}
treats a fetched \texttt{0x00008067} (\texttt{jalr x0,0(x1)}) as ``program finished''
when \texttt{x2} equals \texttt{BASE + MEM\_DEPTH} --- and the register file \emph{resets}
\texttt{x2} to precisely that value. Any benchmark containing a canonical return therefore
terminates early and is graded against \texttt{x15} instead of \texttt{x1}. \texttt{callret}
clears \texttt{x2} at entry to defuse the check; without that, the return-address-stack
workload cannot be expressed at all.

\paragraph{The ecall drain window will overwrite a pass register.} The testbench samples
the register file about five cycles after \texttt{ecall} is fetched. An instruction placed
after the \texttt{ecall} --- a \texttt{fail} block, for instance --- retires inside that
window and can clear the pass marker. Every benchmark now ends with a single \texttt{ecall}
followed only by \texttt{nop}s. This was rediscovered the hard way even though it was
already documented in the suite README, which is the argument for writing it down.

\end{document}
